\documentclass[conference]{IEEEtran}
\IEEEoverridecommandlockouts
\usepackage{cite}
\usepackage{amsmath,amssymb,amsfonts}
\usepackage{algorithm}
\usepackage{algpseudocode}
\usepackage{graphicx}
\usepackage{textcomp}
\usepackage{xcolor}
\usepackage{booktabs}
\usepackage{listings}
\usepackage{url}

\lstset{
  basicstyle=\ttfamily\footnotesize,
  breaklines=true,
  frame=single,
  captionpos=b
}

\begin{document}

\title{AZ-Kerberos: Asynchronous Anomaly-Driven Ticket Policy Enforcement on Stock MIT Kerberos with Cryptographic Decision Binding}

\author{\IEEEauthorblockN{Anonymous Authors}
\IEEEauthorblockA{\textit{Department of Computer Security and Applied AI} \\
Project 22AIE314 Computer Security\\
\texttt{\{contact\}@domain.edu}}
}

\maketitle

\begin{abstract}
Kerberos version 5 (RFC 4120) is the cornerstone authentication protocol for modern enterprise networks and directory services. However, standard Kerberos architectures suffer from an inherent security asymmetry: ticket issuance decisions are static and binary. Once authenticated, users receive unconstrained service tickets---typically valid for 8 to 10 hours---which cannot be revoked prior to expiration. Consequently, adversaries who obtain valid credentials through phishing, password spraying, or offline Kerberoasting enjoy broad lateral movement across the network without intermediate challenge. While intrusion detection systems (IDS) analyze KDC audit logs, existing solutions operate strictly in an \textit{alert-only} posture, creating an unmitigated exposure window between detection and human analyst remediation. Conversely, deploying complex anomaly-detection models directly on the synchronous authentication request path introduces prohibitive latency penalties (12.9~ms per request in our single-call Python benchmark).

In this paper, we present \textbf{AZ-Kerberos}, a closed-loop, risk-adaptive ticket issuance system designed for stock MIT Kerberos without protocol modifications or custom client software. AZ-Kerberos decouples heavy anomaly scoring from synchronous ticket issuance through a three-tier architecture: (1) an asynchronous, log-driven scoring engine that evaluates sliding-window behavioral features via unsupervised Isolation Forests; (2) a lockless, memory-mapped binary risk table utilizing FNV-1a hashing with open-addressing linear probing; and (3) a lightweight C \texttt{kdcpolicy} plugin executing on the KDC request path that enforces adaptive ticket lifetime and renewal caps or request denials in \textbf{about 14--20~nanoseconds}. Furthermore, we identify the \textit{Victim Lock-Out Dilemma} (Threat T1)---wherein attackers intentionally trigger failed authentications to deny legitimate users---and add a multi-signal escalation rule and a realm circuit breaker. Finally, AZ-Kerberos binds every policy decision and table publication into an Ed25519-signed, hash-chained audit log achieving \textbf{64,476 events/second} throughput. In real end-to-end runs against a live MIT krb5 1.20.1 KDC (3 repeats, paced attackers, every outcome read from the KDC's own responses), the loop worked as designed: after detection, newly issued tickets carried the capped lifetime (3,600~s instead of 36,000~s), no victim of a 15-target password spray was capped or denied, every audit log verified, and the plugin failed open when the table was missing, corrupt or stale. The effect on attacks was modest: the scorer flagged a Kerberoasting attacker only after 6 distinct services (6 of 18 AES and 6 of 10 RC4 tickets were issued at full lifetime; 15 of 18 AES tickets in an unpaced burst), the denial level was never reached in a live attack run, and a static threshold rule did about as well (8 of 18 full AES tickets). A lifetime cap does not stop offline cracking of a ticket already issued; only denial withholds it.
\end{abstract}

\begin{IEEEkeywords}
Kerberos Authentication, Network Security, Risk-Adaptive Access Control, Anomaly Detection, Pluggable Policy, Tamper-Evident Logging.
\end{IEEEkeywords}

\section{Introduction}
Kerberos version 5 \cite{rfc4120} serves as the default authentication protocol for contemporary enterprise environments, Active Directory domains, and high-performance computing clusters. Kerberos relies on a trusted Key Distribution Center (KDC) comprising an Authentication Service (AS) for issuing Ticket Granting Tickets (TGTs) and a Ticket Granting Service (TGS) for issuing application-specific service tickets.

Despite its ubiquitous adoption, the protocol's core security model embodies a fundamental assumption of static trust: authentication decisions are atomic, binary, and unchangeable. Once the KDC verifies a client's credentials, it issues tickets valid for fixed, long durations---commonly 10 hours by default---with renewable limits extending up to seven days. Within this temporal window, issued tickets cannot be selectively revoked or constrained by the KDC without re-keying the underlying services or wiping client caches \cite{wiefling2022tops}.

This architectural limitation is aggressively exploited by modern lateral-movement and post-exploitation tradecraft:
\begin{enumerate}
    \item \textbf{Kerberoasting:} An attacker with standard user credentials requests TGS service tickets for target Service Principal Names (SPNs) encrypted under service account keys, taking the tickets offline to perform high-speed dictionary cracking \cite{kotlaba2020icissp}.
    \item \textbf{Password Spraying and AS-REP Roasting:} Attackers systematically probe user accounts or abuse principals with pre-authentication disabled (\texttt{+preauth}) to acquire crackable AS response payloads \cite{grippo2023kerberos}.
    \item \textbf{Stolen Credential Exploitation:} Compromised user credentials permit arbitrary service enumeration and resource access within the active ticket lifetime.
\end{enumerate}

\subsection{The Limitations of Alert-Only Intrusion Detection}
To combat these techniques, existing literature and commercial monitoring solutions rely on host and KDC event log parsing. However, all existing Kerberos detection mechanisms operate in an \textit{alert-only} paradigm. When an anomaly is detected, an alert is queued in a SIEM dashboard. Between the generation of the alert and human intervention (often exceeding tens of minutes or hours), the adversary has already retrieved the necessary Kerberos ticket hashes, moved laterally, or established persistent backdoors.

\subsection{The Dilemma: In-Line Latency vs. Decoupled Adaptation}
The intuitive remedy---evaluating machine learning anomaly detection models directly on the synchronous ticket issuance path---is unattractive. The KDC answers each request inside a single request loop, so any per-request model call adds directly to client-visible latency. In our single-call benchmark (Python/scikit-learn, synthetic input), an Isolation Forest takes \textbf{12.9~ms} per request, which on a busy KDC would queue requests behind the model; we did not benchmark KDC throughput. Conversely, modifying the Kerberos protocol wire format to support progressive re-authentication breaks existing client libraries, operating system Kerberos SSPs, and RFC 4120 conformance.

\subsection{Contributions}
To resolve this tension, we present \textbf{AZ-Kerberos}, a closed-loop adaptive issuance design for stock MIT Kerberos. In our partial literature search (Section~V) we found no prior work that combines these elements. Our specific contributions are as follows:
\begin{itemize}
    \item \textbf{Decoupled Architecture with Nanosecond Enforcement:} We design a decoupled architecture separating asynchronous feature extraction and model scoring from request-path policy enforcement. By utilizing a memory-mapped binary open-addressing table indexed via FNV-1a-64 hashing, the C \texttt{kdcpolicy} module enforces ticket lifetime caps, renewal caps, and denials in \textbf{about 14--20~nanoseconds}.
    \item \textbf{Stock Compatibility and Zero Protocol Changes:} AZ-Kerberos requires zero modifications to the Kerberos wire protocol, ticket formats, or client operating systems. It builds entirely upon the standard MIT Kerberos pluggable policy architecture (\texttt{kdcpolicy} vtable interface introduced in krb5 1.16).
    \item \textbf{Resolution of the Victim Lock-Out Dilemma (Threat T1):} We identify a critical security vulnerability inherent to reactive anomaly detection in identity systems---where an adversary intentionally sprays authentication failures against victim principals to weaponize the detector into an automated DoS tool. We design and implement an asymmetric escalation policy, a two-signal requirement for denial, and a realm-wide circuit breaker. In our 15-target password-spray runs no victim was capped or denied; we do not claim a formal proof, and the denial level was never reached in a live run.
    \item \textbf{Cryptographically Bound Decision Audit:} Every policy adjustment, published risk table hash, and active enforcement event is linked into an Ed25519-signed, hash-chained audit log supporting 64,476 events/second.
    \item \textbf{Full Empirical Validation on Live MIT KDC:} We compile, install, and evaluate AZ-Kerberos on a live MIT Kerberos deployment with 50 user principals and 20 service principals under live Kerberoasting and password-spray workloads, reporting latency overhead, fail-open resilience, audit integrity and attack mitigation (modest, see Section~IV-D).
\end{itemize}

\subsection{Related Work}
\textit{Kerberos detection} work is alert-only: Kotlaba combines machine learning with honeypots (7/7 true positives, 0 false positives) \cite{kotlaba2020icissp} and a one-class SVM (3 true positives, 34 false positives, on 3 positive events) \cite{kotlaba2021icissp}; Grippo and Kholidy inspect Windows events for forged tickets and report no quantitative result \cite{grippo2023kerberos}; MITRE DET0144 suggests alerting on ticket lifetime \cite{mitredet0144}. None of these changes what the KDC issues. \textit{Kerberos extensions} such as the 2025 KTCGM proposal \cite{ktcgm2025} pair a threshold-shared KDC with ML anomaly scoring that triggers step-up multi-factor authentication; that design is not implemented or measured in the text and needs protocol and client changes. MIT's \texttt{kdcpolicy} hook \cite{k5wikipolicy} is the enforcement point we build on, but ships no scoring or adaptation. \textit{Risk-based authentication} targets interactive web sign-in \cite{wiefling2022tops}, and tamper-evident logging \cite{crosby2009usenix} is not tied to access decisions. In our partial literature search we found no work combining asynchronous scoring, issuance changes on unmodified MIT krb5 and a signed hash-chained decision log; this is a weak negative result, as IEEE Xplore, ACM DL and the major security conferences were not searched.

\section{System Architecture and Design}
AZ-Kerberos organizes functionality into three distinct decoupled components operating across different privilege boundaries:
\begin{enumerate}
    \item \textbf{C Policy Module (\texttt{kdcpolicy\_azk.c}):} Hooks directly into \texttt{check\_as} and \texttt{check\_tgs}. When a client requests credentials, the plugin performs an $O(1)$ lockless memory-mapped lookup in the published risk table. If found, it restricts ticket duration via \texttt{lifetime\_out} and \texttt{renew\_lifetime\_out}, or returns \texttt{KRB5KDC\_ERR\_POLICY} to abort issuance.
    \item \textbf{Asynchronous Scoring Pipeline:} An external Python daemon tails KDC syslog/journal events, extracts 6-dimensional sliding-window features ($\Delta t = 300$~s), scores events using a benign-trained Isolation Forest, and maps raw scores to discrete policy levels using empirical quantiles ($\alpha \in \{0.02, 0.005, 0.001\}$).
    \item \textbf{Cryptographically Bound Audit Log:} Records every risk update, table publication (SHA-256 hash), and KDC policy application in a tamper-evident SHA-256 hash chain with Ed25519-signed checkpoints written every $\kappa$ records ($\kappa = 100$ in the reported run). Records after the last checkpoint are hash-linked but unsigned.
\end{enumerate}

\section{Victim Lock-Out Mitigation (Threat T1)}
An external attacker spraying passwords against a victim account causes authentication failures. If failure counts trigger level 3 (denial), the attacker can systematically lock out legitimate employees without possessing credentials. AZ-Kerberos prevents this through:
\begin{enumerate}
    \item \textbf{Preauth Failure Cap:} Pre-authentication failures alone ($f_{\text{distinct\_svc}} = 0 \land f_{\text{rc4}} = 0$) are strictly capped at Level 2 (10-minute lifetime cap, never denial).
    \item \textbf{Two-Signal Requirement for Denial:} Level 3 denial requires at least two independent anomalous signals.
    \item \textbf{Global Circuit Breaker:} If active restrictions exceed 25\% of realm principals, the breaker trips, downgrading all further escalations to Level 1.
\end{enumerate}

\section{Empirical Evaluation}
We evaluated AZ-Kerberos on an experimental testbed running stock MIT Kerberos 1.20.1 in Ubuntu 24.04 LTS.

\subsection{Experiment X4: Request Latency Overhead}
As shown in Table~\ref{tab:microbench}, direct in-memory C lookup took \textbf{14.2~ns} per lookup for 5,000 entries in the reported run (a later re-run on the same machine gave 18.3~ns for 50 entries and 19.6~ns for 5,000, so we quote about 14--20~ns). In live KDC client testing across 60 repeated trials, median end-to-end \texttt{kinit} latency was 467.4~ms (stock KDC), 382.4~ms (empty table), 256.6~ms (50 entries) and 252.5~ms (5,000 entries). The stock KDC is the slowest row, which points to run-order and WSL noise rather than a speed-up, so these runs show only that the plugin adds no overhead measurable above that noise.

\begin{table}[htbp]
\caption{Illustrative lookup latency vs. one in-line model call (C hash lookup vs. Python/sklearn on a synthetic 10-feature input; not a like-for-like comparison)}
\label{tab:microbench}
\centering
\resizebox{\columnwidth}{!}{\begin{tabular}{lrr}
\toprule
\textbf{Model / Architecture} & \textbf{Latency} & \textbf{Relative Speedup} \\
\midrule
In-Line Isolation Forest (Python) & 12,900,000 ns & Baseline ($1.0\times$) \\
In-Line Logistic Regression & 93,500 ns & $138\times$ \\
\textbf{AZ-Kerberos C Table (50 entries)} & \textbf{13.4 ns} & \textbf{962,686}$\times$ \\
\textbf{AZ-Kerberos C Table (5,000 entries)} & \textbf{14.2 ns (re-run 19.6)} & \textbf{$\sim\!10^{6}$}$\times$ \\
\bottomrule
\end{tabular}}
\end{table}

\subsection{Experiment X5: Audit Log Throughput and Tamper Matrix}
The audit logging engine achieved \textbf{64,476.9 events/s} append throughput and \textbf{81,748.2 records/s} verification throughput ($\kappa = 100$). Across 6 adversarial test vectors (bit-flips, record deletion, record reordering, unauthorized insertion, forged keys, and anchored truncation), all 6 were detected (our own test vectors, not an independent evaluation).

\subsection{Experiment X6: Fail-Open Resilience}
Under missing, corrupted, truncated, and stale ($>300$~s) risk tables, AZ-Kerberos consistently fell open to standard 36,000~s (10-hour) ticket issuance without dropping service. Under fresh tables, lifetime caps and denials were enforced strictly as specified.

\subsection{Experiments X1--X3: Live Attack Mitigation}
These are real end-to-end runs on the live MIT krb5 1.20.1 KDC (\texttt{tests/run\_live\_experiment\_real.py}). Each request is a real \texttt{kinit} or \texttt{kvno} call; the engine is fed only from the KDC's own syslog; each outcome (full, capped, denied) is read from the KDC's response (ticket start and end times). The scorer runs as a separate tick every 1~s, so requests between ticks are unprotected. Three modes were compared: unprotected KDC, a static threshold rule (same engine, cap-only), and AZ-Kerberos (Isolation Forest trained on 720 real benign events, two-signal denial rule, signed audit log). Per run: 200 benign requests from 25 users, an 18-service AES Kerberoast, a 10-service RC4 Kerberoast, and a spray of 3 wrong passwords against each of 15 victims who then log in legitimately. The attacker is unpaced or paced at 1 or 3~s per request; 3 repeats with 3 benign sequences. Table~\ref{tab:attack_mitigation} gives means. Capped tickets are issued with a 3,600~s lifetime instead of 36,000~s. A cap limits later replay but not offline cracking of a ticket already issued; only denial withholds a ticket, and denial was \emph{never} triggered in a live attack run (its rule needs at least 10 distinct services and 5 RC4 requests, which neither attack produced; we did not tune it after seeing the results). AZ-Kerberos restricted the AES attacker from its 7th request (6 full tickets) versus the 9th (8 full) for the threshold rule; RC4 results were equal. An unpaced burst finished before the 1~s tick in most runs: AZ-Kerberos restricted it in 1 of 3 runs (15 of 18 tickets full on average), the threshold rule in 0 of 3. A 0.25~s tick (2 repeats) gave 13.5 full AES tickets for the burst. No spray victim was capped or denied under AZ-Kerberos; the threshold rule capped 10 of 15 only because all lab clients share 127.0.0.1, so its per-source spray feature pools all actors. Benign false restrictions: 4 of 600 requests (one of 25 users in one of the 3 benign sequences); 0 of 400 with the 0.25~s tick. All AZ-Kerberos audit logs (13 of 13 runs) verified. Stolen-credential use from a new host was not run, and an earlier synthetic X1--X3 run was discarded.

\begin{table}[htbp]
\caption{Real live-KDC results, mean of 3 repeats, tick 1~s (full / capped / denied tickets)}
\label{tab:attack_mitigation}
\centering
\resizebox{\columnwidth}{!}{\begin{tabular}{lrrr}
\toprule
\textbf{Metric} & \textbf{Unprotected} & \textbf{Threshold} & \textbf{AZ-Kerberos} \\
\midrule
AES (18 req.), paced 1 s / 3 s & 18 / 0 / 0 & 8 / 10 / 0 & 6 / 12 / 0 \\
AES (18 req.), unpaced burst & 18 / 0 / 0 & 18 / 0 / 0 & 15 / 3 / 0 \\
RC4 (10 req.), paced 1 s / 3 s & 10 / 0 / 0 & 6 / 4 / 0 & 6 / 4 / 0 \\
RC4 (10 req.), unpaced burst & 10 / 0 / 0 & 10 / 0 / 0 & 10 / 0 / 0 \\
Spray victims (15) logging in & 15 / 0 / 0 & 5 / 10 / 0 & 15 / 0 / 0 \\
Benign requests restricted (of 200) & 0 & 0 & 1.33 \\
Time to first restriction, AES (1 s / 3 s) & -- & 8.6 / 24.5 s & 6.4 / 18.5 s \\
\bottomrule
\end{tabular}}
\end{table}

\section{Limitations}
\begin{itemize}
\item \textbf{Literature coverage is partial.} About 20 sources were read per pillar; IEEE Xplore, ACM DL, NDSS, S\&P and CCS were not searched directly. ``No prior work feeds a risk score into Kerberos ticket issuance'' is a weak negative result, not a proof. Qomplx patent US12580898 was read: it detects forged tickets and destroys them after the fact using a ledger of authentication objects, with no lifetime adaptation, no signed hash chain and no stock-krb5 deployment (partial overlap only); Qomplx US11968227, US12438906 and IEEE Computer 2024 (DOI 10.1109/MC.2024.3434535) are unread. Post-quantum PKINIT on stock MIT krb5 already exists (kurbu5-pkinit), so we make no post-quantum claim.
\item \textbf{Evaluation scale.} X1--X3 are scripted attacks on a lab realm of 50 users: 3 repeats per cell, 3 benign sequences, one scorer configuration, one threshold setting, no statistical testing, all clients on 127.0.0.1 (so IP features cannot discriminate), stolen-credential use from a new host not run. Latency percentiles use 60 samples per configuration.
\item \textbf{Detection is late and denial never fired.} The scorer restricts only after about 6 distinct services and an unpaced burst finishes before the tick; capped tickets remain crackable; the denial level was not reached in any live attack run. The Isolation Forest showed no demonstrated advantage over the hand-set threshold rule (6 versus 8 full AES tickets, equal RC4).
\item \textbf{Client IP.} The plugin receives principals and request flags, not the source address, so enforcement is keyed by principal and IP correlation lives in the asynchronous engine.
\item \textbf{Low-and-slow attacks.} One service request every few days stays below the sliding-window features.
\end{itemize}

\section{Conclusion}
AZ-Kerberos shows that a risk score computed off the request path can reach the KDC's ticket issuance on stock MIT Kerberos 1.20.1 without protocol changes: real tickets issued after detection carry the capped lifetime, the in-process lookup takes about 14--20~ns, and the signed, hash-chained audit log verified in every run. It does not yet stop Kerberoasting: the attacker kept 6 of 18 AES tickets at full lifetime (15 of 18 in a burst), capped tickets stay crackable, denial was never reached, and a threshold rule did about as well. Shortening the detection delay and testing denial against real attack tooling are the main open problems. Results come from a single lab realm and should be read with the limitations above.

\begin{thebibliography}{00}
\bibitem{rfc4120} C. Neuman et al., ``The Kerberos Network Authentication Service (V5),'' RFC 4120, IETF, 2005.
\bibitem{wiefling2022tops} S. Wiefling et al., ``What's in a Score? Evaluating the Reliability and User Impact of Risk-Based Authentication,'' \textit{ACM TOPS}, vol. 25, no. 4, 2022.
\bibitem{kotlaba2020icissp} M. Kotlaba, ``Detection of Kerberoasting Attacks Using Machine Learning and Honeypots,'' in \textit{Proc. ICISSP}, 2020.
\bibitem{kotlaba2021icissp} M. Kotlaba, ``An Anomaly-based Intrusion Detection System for Kerberos using One-Class Support Vector Machines,'' in \textit{Proc. ICISSP}, 2021.
\bibitem{grippo2023kerberos} M. Grippo and H. A. Kholidy, ``Detecting Forged Kerberos Tickets in an Active Directory Environment,'' \textit{arXiv:2301.00044}, 2022.
\bibitem{ktcgm2025} R. Almatarneh et al., ``KTCGM: Towards a Novel Solution for Enhancing Kerberos-5 with Threshold Cryptography and ML-Based Anomaly Detection,'' \textit{Int. J. Innovative Research and Scientific Studies}, vol. 8, no. 3, pp. 3646--3662, 2025.
\bibitem{k5wikipolicy} MIT Kerberos Project, ``Projects/KDC TGS Policy plugin,'' K5Wiki, \url{https://k5wiki.kerberos.org/wiki/Projects/KDC_TGS_Policy_plugin}, accessed 2026-11-22.
\bibitem{mitredet0144} MITRE, ``Detect Forged Kerberos Golden Tickets (T1558.001), Detection Strategy DET0144,'' MITRE ATT\&CK, \url{https://attack.mitre.org/detectionstrategies/DET0144/}, accessed 2026-11-22.
\bibitem{crosby2009usenix} S. A. Crosby and D. S. Wallach, ``Efficient Data Structures for Tamper-Evident Logging,'' in \textit{Proc. USENIX Security}, 2009.
\end{thebibliography}

\end{document}
